\textbf{1. Kiến thức cần nhớ}
\begin{itemize}
\item \textbf{Dữ liệu}: Thông tin thu thập được bao gồm chữ, số, hình ảnh, âm thanh\dots
\item \textbf{Phân loại dữ liệu}:
\begin{itemize}
\item \textbf{Dữ liệu định tính} (dữ liệu phi số): Biểu thị bằng chữ, kí hiệu.
\begin{itemize}
\item Dữ liệu định tính \textit{không thể sắp thứ tự}: ví dụ màu sắc yêu thích (xanh, đỏ, vàng), quốc tịch\dots
\item Dữ liệu định tính \textit{có thể sắp thứ tự}: ví dụ xếp loại học lực (xuất sắc, giỏi, khá, trung bình), kích cỡ quần áo (S, M, L, XL)\dots
\end{itemize}
\item \textbf{Dữ liệu định lượng} (dữ liệu số): Biểu thị bằng các số thực.
\begin{itemize}
\item \textit{Dữ liệu rời rạc}: Nhận các giá trị số nguyên rời rạc, thường thu được từ phép đếm (ví dụ: số học sinh, số bàn thắng, số lỗi chính tả\dots).
\item \textit{Dữ liệu liên tục}: Có thể nhận mọi giá trị trong một khoảng số thực, thường thu được từ phép đo (ví dụ: chiều cao, cân nặng, nhiệt độ, thời gian chạy\dots).
\end{itemize}
\end{itemize}
\item \textbf{Tính hợp lí của dữ liệu}: Dữ liệu phải đảm bảo tính chính xác, nằm trong khoảng cho phép của thực tế (ví dụ: điểm kiểm tra không thể âm hoặc lớn hơn $, tuổi học sinh lớp 8 không thể là $ hay số âm).
\end{itemize}

\textbf{2. Phương pháp giải các dạng}

%% DANG 8D5?18-1
\textbf{Dạng 1. Phân loại dữ liệu (định tính, định lượng; rời rạc, liên tục)}
\begin{enumerate}
\item Xem xét giá trị dữ liệu là chữ hay số:
\begin{itemize}
\item Nếu là chữ/kí hiệu $\rightarrow$ Dữ liệu định tính (xét tiếp có xếp thứ bậc được không).
\item Nếu là số $\rightarrow$ Dữ liệu định lượng.
\end{itemize}
\item Với dữ liệu định lượng:
\begin{itemize}
\item Thu thập bằng cách đếm các đơn vị nguyên vẹn $\rightarrow$ Dữ liệu rời rạc.
\item Thu thập bằng dụng cụ đo lường (thước, cân, đồng hồ) $\rightarrow$ Dữ liệu liên tục.
\end{itemize}
\end{enumerate}

%% DANG 8D5?18-2
\textbf{Dạng 2. Thu thập dữ liệu, nhận biết dữ liệu không hợp lí}
\begin{enumerate}
\item Các phương pháp thu thập: quan sát trực tiếp, làm thí nghiệm, lập phiếu khảo sát (phỏng vấn), thu thập từ nguồn có sẵn (sách, báo, mạng Internet uy tín).
\item Tiêu chí kiểm tra tính hợp lí:
\begin{itemize}
\item Đúng kiểu dữ liệu (ví dụ: điểm số phải là số, không thể là chữ).
\item Nằm trong phạm vi cho phép của thực tế (tỉ lệ phần trăm không vượt quá \%$, tổng các tỉ lệ thành phần trong cùng mẫu phải bằng \%$).
\end{itemize}
\end{enumerate}
